\documentclass[12pt]{article}
\usepackage[margin=1in]{geometry}
\usepackage{amsmath,amssymb,booktabs,enumitem,hyperref}
\title{Investigation Log: Zero Forcing Numbers of $P(n,k)$ for $k\ge4$}
\author{Aryan Padarthi}
\date{September 23, 2026}
\setlength{\emergencystretch}{3em}
\begin{document}
\maketitle

\section*{Entry 1: How this problem was chosen}
\textbf{Date:} September 23, 2026\\
This is the ninth research target attempted in this project. The previous eight
were abandoned, and the reason matters: in every case the problem turned out to
be already solved, and in every case I had inferred that it was open from the
\emph{absence} of a published result rather than from a positive statement.

\begin{itemize}[leftmargin=*]
\item Four attempts on Turaev genus and signature bounds for positive braid
closures --- an area with an active professional research group.
\item Periodicity of the octal game \texttt{0.45}. I derived a real structural
mechanism, then found the Guy--Smith periodicity theorem (1956), which settles
it outright once one has computed past $n=1018$. I had verified to $60{,}000$.
\item A ``rare-index invariant'' I derived for octal games, which is exactly
Berlekamp's Sparse Space Phenomenon, described in \emph{Winning Ways} c.\ 1980.
\item Persistent homology of linkage configuration spaces. Farber--Sch\"utz
(2007) give closed Betti-number formulas for planar polygon spaces, and
Ba--Zhou (arXiv:2506.06781, 2025) prove the self-avoiding moduli spaces are
diffeomorphic to Euclidean space --- i.e.\ contractible, so persistent homology
would correctly compute nothing at all.
\end{itemize}

The corrected procedure: before any investment, ask \emph{what standard result
would make this trivial}, and require a positive statement of openness. This
target passes that test. Krishnan (arXiv:2607.19412, July 2026) closes with:
``Determining $Z(P(n,k))$ for general $k$ --- the stabilization threshold and a
rigorous lower bound for $k\ge4$ --- remains open.'' That is an explicit
assertion in a paper two months old, not an inference from silence.

\section*{Entry 2: A validated engine}
\textbf{Date:} September 23, 2026\\
No value was used for research before the engine reproduced every published
value. The validation set:
\begin{itemize}[leftmargin=*]
\item Krishnan's Table 1, $Z(P(n,3))$ for $7\le n\le20$:
$6\,6\,6\,8\,7\,7\,8\,8\,8\,8\,8\,8\,8\,8$ --- reproduced exactly, including
the non-monotone spike at $n=10$ and the corrected $Z(P(12,3))=7$ that
falsifies \cite{RST}'s Theorem 3.6.
\item $Z(P(n,2))=6$ for $10\le n\le16$; $Z(P(2k+1,k))=6$ for $k=5,6,7$.
\item Independent cross-check against brute force on small cases.
\end{itemize}
The engine independently found the witness $\{u_0,\dots,u_6\}$ for
$Z(P(12,3))=7$ --- precisely the configuration Krishnan identifies as the one
the published proof never tested.

\section*{Entry 3: Three approaches that failed}
\textbf{Date:} September 23, 2026\\
Recorded because the negative results were informative and cost real time.

\textbf{(a) Fort-based ILP (scipy).} Exact in principle: a set is zero forcing
iff it meets every fort, so $Z$ is a minimum hitting set. With a partial fort
family the ILP is a relaxation, giving a valid lower bound, and any returned
set is verified independently by running the forcing rule --- so the method
cannot report a wrong value, only fail to converge. It failed to converge:
\texttt{scipy.optimize.milp} has no lazy-constraint callback, so every round
re-solves from scratch. Stalled past $P(11,3)$, which brute force does
instantly.

\textbf{(b) CP-SAT, time-indexed formulation.} Correct on every instance it
finished, all witnesses verified. Also far too slow --- $30$s on $P(13,3)$,
stuck on $P(15,3)$. The obstruction is structural, not incidental: finding a
forcing set is easy, \emph{proving minimality} is hard, and the formulation's
LP relaxation is too weak to prune. The exponential cost moves from enumeration
into branch-and-bound rather than disappearing.

\textbf{(c) Disjoint forts as a lower-bound certificate.} If $F_1,\dots,F_t$
are pairwise disjoint forts then $Z\ge t$, with no case analysis --- which is
exactly the kind of argument the open lower bound needs. This fails for an
information-theoretic reason: forts in $P(n,k)$ have $7$--$13$ vertices while
the graph has only $2n\le56$, so a greedy search found at most \emph{two}
disjoint forts in every case, giving $Z\ge2$. Not a tuning problem; the
certificate cannot exist at these sizes.

What survived was the simplest thing: exhaustive enumeration in C with
$128$-bit masks, plus a rotation-symmetry reduction, plus threads.

\section*{Entry 4: Results}
\textbf{Date:} September 23, 2026\\
\textbf{Proved.} For all $k\ge2$ and $n\ge2k+3$, $Z(P(n,k))\le2k+2$, by an
explicit forcing set $\{u_0,\dots,u_{2k+1}\}$ and a rotation-invariance
bootstrap. The bound itself is \cite{RST}'s Theorem 2.2; the contribution is a
proof uniform in $k$ (Krishnan's clean replacement covers only $k=3$) in a
setting where the original paper is known to contain a false theorem. Every
step machine-verified for $2\le k\le8$.

\textbf{Computed.} Exact values of $Z(P(n,k))$ for $k=4$ ($9\le n\le28$),
$k=5$ ($11\le n\le25$), $k=6$ ($13\le n\le26$). These are new; the literature
has no exact values for $k\ge4$.

\textbf{Conjectured, not proved.} That $Z(P(n,k))=2k+2$ for large $n$, with
thresholds $10,13,18,23$ for $k=2,3,4,5$. I decline to fit a formula to four
points; $4k+2$ fails at $k=3$ and $k=5$.

\textbf{Open.} The matching lower bound. Nothing here bears on it.

\textbf{Unresolved in our own data.} Whether $Z(P(n,6))$ ever attains $14$.
The maximum observed for $n\le26$ is $12$, and values with $\gcd(n,6)=1$ appear
to cap lower than those with $\gcd(n,6)>1$. Since the inner vertices form
$\gcd(n,k)$ disjoint cycles this is structurally plausible, but with the
ceiling unattained anywhere in range it is not claimed as a finding.

\section*{Entry 5: A correction I had to make to my own claim}
\textbf{Date:} September 23, 2026\\
After seeing $Z(P(n,4))=10$ for $n=18,\dots,22$ --- five consecutive values ---
I concluded the sequence had stabilized. The $k=5$ row then showed
$Z=10$ at $n=17,\dots,21$, also five consecutive values, followed by $11$ at
$n=22$ and $12$ at $n=23$. Five repeats are not evidence of stabilization. The
$k=4$ claim survives only because $2k+2$ is a \emph{proven} ceiling there, so
the sequence cannot climb --- but it dropped at $n=12$ and $n=16$, so it could
in principle drop again. This is the same warning Krishnan gives about
extrapolating from finite computation, arriving on my own data within an hour
of my reading it.


\section*{Entry 6: A strategic reassessment}
\textbf{Date:} September 23, 2026\\
Asked to assess whether the work so far constituted a project, the honest
answer was no. Stripped of presentation, the content was: new computed values,
plus a \emph{rederivation} of somebody else's theorem. The open problem --- the
lower bound --- was untouched. Compared against the project that qualified from
this region last year (an IFS construction on the Riemann $\xi$-function, which
proves existence, uniqueness, an antipodal symmetry, and a dimension bound),
one rederived theorem and a table is thin.

Three routes to original content were weighed:
\begin{itemize}
\item \emph{The lower bound.} Highest value, lowest odds. It had already
defeated a published proof, an automated prover (a logged attempt of
2026-09-16, closed without progress), and our own disjoint-fort attempt.
\item \emph{Exceptional values.} Tempting, since upper bounds are constructive.
Abandoned after checking: the drops are transient. For $k=4$, the values at
$n\equiv0\pmod4$ run $6,8,10,10,10$ --- they climb to the ceiling like
everything else, so there is no permanent exceptional family to characterise.
\item \emph{Generalise the method, not the family.} Chosen. The rotation
bootstrap is not specific to $P(n,k)$; it needs only a cyclic automorphism.
\end{itemize}
The third worked, and it is what turned the day's output from a table into
theorems.

\section*{Entry 7: Three lower-bound techniques, all insufficient}
\textbf{Date:} September 23, 2026\\
Before abandoning lower bounds, each standard route was tested to destruction
rather than assumed hopeless.
\begin{itemize}
\item \textbf{Disjoint forts.} $t$ pairwise disjoint forts give $Z\ge t$ with
no case analysis. Greedy search found at most \textbf{two} in every instance.
Forts here have $7$--$13$ vertices and the graph has $2n\le56$; packing $2k+2$
of them is not a tuning problem but an impossibility.
\item \textbf{The fort LP.} Computed exactly by cutting planes, separating with
a minimum-weight-fort ILP each round. $\mathrm{LP}/Z\approx0.43$--$0.50$
throughout, worst at $P(18,4)$: $3.21$ against $Z=10$. An integrality gap of
about $2$, widening with $k$.
\item \textbf{Hoffman-type spectral bound.} Gives $1.2$--$2.9$ against true $Z$
of $6$--$12$; also degrades in $k$.
\end{itemize}
This is a negative result worth stating: the lower bound resists every
general-purpose technique available, which is why it is open. It also
retroactively explains the earlier solver failures --- branch and bound was
being asked to close a factor-of-two gap by branching.

\section*{Entry 8: Two theorems}
\textbf{Date:} September 23, 2026\\
The bootstrap was isolated into two elementary lemmas (closure commutes with
automorphisms; a filled outer set forces the inner set) and then applied.

\textbf{$Z(I(n,j,k))\le 2(j+k)$ for $n\ge2(j+k)+1$.} The chain is: $u_i$ forces
$v_i$ for $j\le i\le m-1-j$; then $v_{j+k}$ forces $v_{j+2k}$; then $u_{j+2k}$
forces $u_m$. Taking $m=2(j+k)$ makes the second step work because
$m-j-2k=j$ is exactly the first index filled in step 1. Verified in 106 cases.
Setting $j=1$ recovers $Z(P(n,k))\le2k+2$, so the earlier rederivation is now
a corollary rather than a separate result.

\textbf{$Z(DP(n,k))\le 4k+4$ for $n\ge2k+3$.} The same chain, doubled and run
in parallel on both outer cycles. Verified for $1\le k\le5$. No zero forcing
results for this family were found in the literature. Our computed values give
$Z(DP(n,1))=8=4\cdot1+4$ for $6\le n\le12$, and subsequently
$Z(DP(10,2))=12=4\cdot2+4$, so the bound is attained at $k=1$ and $k=2$ --- it
is sharp there, not merely an upper estimate.

\section*{Entry 9: A prediction, recorded before the test}
\textbf{Date:} September 23, 2026\\
The thresholds at which $Z(P(n,k))$ first attains $2k+2$ are $10,13,18,23$ for
$k=2,3,4,5$. For $k\ge3$ these are exactly $5k-2$; $k=2$ is the exception.
The formula predicts $n_0(6)=28$.

At the time of writing the $k=6$ computation has reached $Z(P(27,6))=13$, which
is consistent (the ceiling $14$ is not yet attained) but does not test the
prediction. $P(28,6)$ is running. This is recorded now, before the answer is
known, so that a confirmation counts as a test rather than a fit --- the same
discipline that Entry 5 failed to apply.


\section*{Entry 10: The prediction was tested, and the conjecture died}
\textbf{Date:} September 24, 2026\\
Entry 9 recorded, before the computation finished, that $n_0(6)=5\cdot6-2=28$
and hence $Z(P(28,6))=14$. The run took just over fourteen hours.

\textbf{$Z(P(28,6))=14$}, with witness $\{u_0,\dots,u_{13}\}$. The narrow
prediction was correct: $28$ is the least $n$ at which $k=6$ attains its
ceiling, exactly as $5k-2$ says.

\textbf{$Z(P(29,6))=12$.} And this refutes the conjecture. It was stated as
``$Z(P(n,k))=2k+2$ for all $n\ge n_0(k)$'' --- a \emph{stabilization} claim ---
and $29>28$ with $Z=12\ne14$. So the prediction and the conjecture came apart:
the thing I could check quickly was right, and the thing it was evidence for
was wrong.

The cause is $\gcd(n,k)$, which had been visible in the $k=6$ data for some
time and which I twice set aside because $k=4$ and $k=5$ looked uniform.
Sorting by $d=\gcd(n,6)$, each class climbs on its own schedule: $d=2$ reaches
$14$ at $n=28$ while the coprime class has only reached $12$ by $n=29$. At
$k=4$ every class had already converged inside the computed range, so no
separation was visible; at $k=6$ the range finally became wide enough to see
it. The apparent thresholds $10,13,18,23$ are where the slowest class happened
to arrive, not values of a formula.

Two things worth recording about the process rather than the result. First,
the prediction was useful precisely because it was written down in advance ---
had I fitted $5k-2$ after seeing $Z(P(28,6))=14$, I would have called it a
confirmation and never computed $n=29$. Second, this is the third time in two
days that a pattern read off a short run has failed on the next data point
(five repeats at $k=5$; ``$Z$ tracks $n$'' for I-graphs; now $5k-2$). The
consistent lesson is that in this family, sequences look stable well before
they are.


\section*{Entry 11: A second conjecture, refuted by data I already had}
\textbf{Date:} September 25, 2026\\
From $20$ computed I-graphs, all satisfying $Z\le2(d_1+d_2)$ with
$d_1=\gcd(n,j)$, $d_2=\gcd(n,k)$, I wrote a conjecture into the paper
asserting that bound in general. It is false, and the counterexamples were
already in the validation table.

Every one of those $20$ graphs had $d_1,d_2>1$ --- that was the selection
criterion for ``genuinely new'' I-graphs. Setting $j=1$ gives $d_1=1$, and with
$\gcd(n,k)=1$ as well the bound reads $Z\le4$, while $Z(P(n,k))=2k+2$ grows
without bound. $Z(P(11,3))=7$ and $Z(P(17,4))=9$ refute it immediately. I had
computed both values days earlier.

The error was generalising from a sample selected by a condition, to a
statement about everything, without testing the boundary the condition
excluded. The previous conjecture ($n_0=5k-2$) at least took a fourteen-hour
computation to kill; this one needed a table lookup.

What survives is narrower and still useful. The observation holds as stated for
$d_1,d_2>1$, and a separate computation supports the underlying intuition: the
minimum \emph{bootstrapping} set for $I(18,3,8)$ --- the object a
uniform-in-$n$ proof actually needs --- has size $10=2(d_1+d_2)$ against
$2(j+k)=22$, and decomposes as three consecutive vertices in each of the three
outer cycles plus one inner vertex. Three consecutive vertices in an outer
cycle is precisely what lets the middle one force its spoke. So the route to
beating the $Z(I(n,j,k))\le 2(j+k)$ bound is a proof adapted to the cycle
decomposition
rather than to consecutive indices --- but the target it proves must retain
some dependence on $j$ and $k$, not on $d_1,d_2$ alone.


\section*{Entry 12: A lower bound, at last}
\textbf{Date:} September 25, 2026\\
Three combinatorial routes to a lower bound had failed (disjoint forts, the
fort LP, a spectral bound), each degrading as $k$ grows. The route that works
is not combinatorial at all: $Z(G)\ge M(G)$, the maximum nullity over symmetric
matrices with $G$'s pattern, so a single explicit matrix \emph{proves} a bound.
$P(n,k)$'s rotation symmetry makes such matrices computable.

Three stages, each with a lesson.

\textbf{$d=1$.} Taking $A$ equivariant under $\rho$ splits it into $n$ blocks of
size two, and the singularity condition is linear and homogeneous in four
coefficients --- a nullspace computation, not a search, so the blocks are
singular by construction. This gave $Z(P(12,2))=Z(P(14,2))=6$ in both
directions, the project's first lower bound.

\textbf{$d=2$, and a numerical trap.} Relaxing to $\rho^{2}$ gives ten
parameters and $4\times4$ blocks. My first attempt solved for them numerically
and reported nullity $7$ for $P(14,2)$, where $Z=6$. Since $Z\ge M\ge$ nullity
that is impossible, so the theory caught the bug: the optimiser had shrunk a
required-nonzero edge weight to $3\times10^{-4}$ against a matrix scale of
$4.5$, numerically deleting an edge while still passing the zero-pattern test.
I now verify by block determinants, which are zero by construction, rather than
by thresholding eigenvalues.

Doing $d=2$ algebraically instead revealed a real obstruction: six coefficients
suggest five fitted points, but fitting five always forced $c_r^{2}\le0$ ---
the spoke weight vanishing, decoupling the two halves. Fitting \emph{four}
leaves a one-parameter pencil containing realisable members. That gave $8$ for
$P(18,4)$ and $10$ for $P(24,6)$, and unlike $d=1$ the bound now grows with
$k$.

\textbf{Uniformity, via Chebyshev.} The certificates were still found one $n$ at
a time. But $s=2\cos\theta$ and $t=2\cos k\theta$ means $t=2T_k(s/2)$ exactly,
so the condition is a polynomial in $s$ alone of degree $k+1$ and the singular
blocks are its roots. For $k=2$ the cubic's coefficients are free, so any three
$s$-values can be prescribed --- a construction valid for every $n$, and $6$ is
maximal because a cubic has three roots. This also explains the flat cap I had
observed: for $k\ge2$ the $d=1$ polynomial is forced into a shape with only
three free ratios.

The same substitution for $d=2$ gives degree $k+1$ in $x$, ceiling $2k+2$ ---
exactly the conjectured $Z$, with the shortfall due entirely to realisability.
Carrying it out for $k=4$ yields nullity $8$ for $n=18$ and every even
$22\le n\le40$, so $8\le Z(P(n,4))\le10$ uniformly. No lower bound on
$Z(P(n,k))$ for $k\ge3$ appears in the literature we have seen.

What I take from this: the three failures were all of one kind --- asking a
combinatorial relaxation to certify a quantity it cannot see --- and the fix
was to change the category of argument rather than to work harder inside it.

\section*{Entry 13: The ceiling, and that it is exactly the upper bound}
\textbf{Date:} September 25--26, 2026 \\
\textbf{Command:} {\small\texttt{python verification/recurrence\_order.py}}

I had been describing the certificate method as ``capped near $6$--$8$'' without
knowing what the cap actually was. It has an exact value, and finding it
changed the project.

Every spoke weight $c_i$ is nonzero, so the kernel equation at $u_i$ solves for
$y_i$ as a $3$-term expression in $x$. Substituting into the equation at $v_i$
leaves one scalar recurrence in $x$ whose support is the nine indices
$\{i{-}k{-}1,i{-}k,i{-}k{+}1\}\cup\{i{-}1,i,i{+}1\}\cup\{i{+}k{-}1,i{+}k,i{+}k{+}1\}$,
with both end coefficients nonzero. So the recurrence has order exactly
$2k+2$, and $\operatorname{null}A=\dim\ker(T-I)$ for its monodromy $T$: at most
$2k+2$, with equality iff $T=I$.

Two things fell out that I did not expect. First, the argument never uses
symmetry --- the bound $\operatorname{null}A\le Z$ needs only a nonzero entry in
each direction along each edge --- so the whole combinatorially symmetric class,
$8n$ parameters instead of $5n$, is available. Second, and this is the part
worth remembering: the order $2k+2$ is the statement that no kernel vector
vanishes on $u_0,\dots,u_{2k+1}$, and that set is exactly the forcing set the
rotation bootstrap produces in Theorem~1. \emph{The ceiling of the certificate
method is the upper bound it is trying to match.} There is no slack anywhere;
either the method proves $Z=2k+2$ or nothing does.

Checked before use, as usual: $189$ random matrices over $63$ pairs $(n,k)$ for
the sparsity, the end-coefficient formulas and
$\operatorname{null}A=\dim\ker(T-I)\le2k+2$; $60$ more for
$\det T=(\prod b')(\prod e')/((\prod b)(\prod e))$, for the product form
$\operatorname{null}A=\dim\ker(SU-I)$, and for the forcing-set statement.

\section*{Entry 14: A claim of mine, refuted --- and the refutation is the result}
\textbf{Date:} September 26, 2026 \quad
\textbf{Command:} \texttt{python verification/d1\_ceiling.py}

Entry 12 asserted that the period-one construction ``never exceeds $6$ for
every $k$'', because the symbol
$F(s)=(s+a)L_k(s)+\alpha s+\beta$ has only three free parameters. I wrote a
script to check that assertion exhaustively over all triples of grid roots,
expecting it to confirm the claim. It printed $12$.

The premise was right and the inference was wrong. At a root the condition is
linear, so three roots determine $(a,\alpha,\beta)$ and no fourth can be
imposed --- but prescribing three roots \emph{determines the other $k-2$}, and
those can land on the grid $\Gamma_n=\{2\cos(2\pi m/n)\}$ by themselves. They
do so exactly when arithmetic makes them: a symbol with rational coefficients
has Galois-closed roots, and the Galois orbit of $2\cos(2\pi m/n)$ stays inside
$\Gamma_n$. One rational $F$ places all $k+1$ roots at once.

I had reasoned about how many roots I could \emph{choose} and concluded
something about how many roots could \emph{be there}. The period-two machinery
of Entry 12, built to escape a cap that was not real, is now redundant.

\section*{Entry 15: Cyclotomic certificates, and the first exact values for
$k\ge4$}
\textbf{Date:} September 26, 2026 \\
\textbf{Commands:} {\small\texttt{cyclotomic\_certificates.py},}\\
{\small\texttt{period1\_optimum.py}, \texttt{certified\_lower\_bounds.py}}

Taking the symbol to be a product of minimal polynomials
$\Psi_d$ of $2\cos(2\pi/d)$, with $\sum\deg\Psi_d=k+1$, and imposing the $k-2$
integer conditions for membership in the family, gives certificates uniform in
$n$:
\[
\begin{array}{ll}
 Z(P(n,4))=M(P(n,4))=10 & \text{for every } n \text{ divisible by } 60,\,70
   \text{ or } 90,\\
 Z(P(n,5))=M(P(n,5))=12 & \text{for every } n \text{ divisible by } 24 .
\end{array}
\]
These are the first exact values of $Z(P(n,k))$ for $k\ge4$; \cite{K} records
that range as open, and a week ago I had $8\le Z(P(n,4))\le10$ and nothing at
all for $k=5$. Every certificate has $a,\alpha,c\in\mathbb Z$, so it is an
integer matrix and its nullity is confirmed by exact rational elimination.

The $k=5$ case is the one I like. There $a=d=0$ and $c=1$: the certificate
\emph{is the adjacency matrix} of $P(n,5)$, the symbol is
$sL_5-1=(s-1)(s+1)(s^4-4s^2+1)=\Psi_6\Psi_3\Psi_{24}$, and the whole statement
reduces to counting $m$ with
$2\cos\frac{2\pi\cdot6m}{n}+2\cos\frac{2\pi\cdot4m}{n}=1$. The simplest matrix
in the class was sufficient all along; I had spent weeks on elaborate ones.

Beyond the rational construction, the exhaustive triple scan over $n\le130$
gives the period-one optimum for each $(n,k)$; those certificates were
re-verified in $\mathbb Q(\zeta_n)$, with $2\cos(2\pi m/n)$ represented as
$x^m+x^{n-m}$ modulo $\Phi_n$, so the membership identity, the reality of the
parameters and $c^2\ne0$ are decided exactly. It adds $Z(P(n,3))=8$ for
$n=40,60,80,120$ (a matrix proof of values \cite{K} gets combinatorially) and
$Z(P(n,2))=6$ for $9\le n\le130$, $n\ne10$. It also shows the honest limit:
period one reaches $2k+2$ for $k\le5$ and, over this range of $n$, gets only to
$8$, $10$, $8$ for $k=6,7,8$ against ceilings $14$, $16$, $18$.

\section*{Entry 16: Two numerical traps, both caught by the ceiling}
\textbf{Date:} September 26, 2026

Counting grid roots by $|F(r)|<10^{-8}$ reported nullity $7$ for $k=2$ at
$n=126,128,130$. That is impossible --- the ceiling is $2k+2=6$ --- and the
cause was $\Gamma_n$ becoming dense enough that near-roots passed the test.
Separately, a candidate at $(n,k)=(48,7)$ with $|a|$ enormous against the edge
weights made a threshold relative to $\max|A|$ report nullity $48$.

Both were caught the same way the period-two trap was caught in Entry 12: by an
inequality that must hold. Then it was $Z\ge M$; now it is
the ceiling of Entry 13. It is worth stating as a working rule:
\emph{every numerical search should be wrapped in a bound it cannot legally
violate}, because that bound is the only thing that will tell you the answer is
wrong when it looks right. The fixes were to test each Fourier block against
its own scale and to certify in $\mathbb Q(\zeta_n)$ rather than in floating
point.

A third mistake, of a different kind: my first search of the combinatorially
symmetric class reported a \emph{smaller} maximum nullity on $P(12,2)$ than the
symmetric search did. Since $S(G)$ sits inside that class, this cannot happen,
so the search --- not the mathematics --- was at fault. Seeding the larger search
with the symmetric certificate fixed it. Before that check I had drawn a
conclusion from those numbers, namely that even $8n$ parameters could not reach
$2k+2$ at $k=3$; the conclusion may still be true but the evidence I had for it
was worthless.

\section*{Entry 17: The period-$d$ parameter count is $2d+1$, not $5d-1$}
\textbf{Date:} September 26, 2026 \\
\textbf{Command:} {\small\texttt{python verification/period\_d\_rank.py}}

Entry 15 ended by naming period-$d$ symbols as the fix for the shortage of
parameters, and I wrote in the paper that a period-$d$ symbol ``carries $5d-1$
free ratios, so $d\ge(k+2)/5$ restores the count.'' That was a guess dressed as
a count, and it is wrong.

The honest way to ask the question is: what is the dimension of the set of
symbols actually \emph{reachable}? Most of the $5d$ weights are gauge ---
conjugating by a positive diagonal matrix preserves both the pattern and the
nullity --- so counting weights counts nothing. I computed the rank of the
Jacobian of
\[
  (\text{5$d$ weights})\ \longmapsto\ (\text{monic coefficients of } G(\sigma))
\]
at random points, for $3\le k\le13$ and $1\le d\le7$. In every one of those
cases,
\[
  \operatorname{rank}=\min(2d+1,\;k+1).
\]
At $d=1$ that returns $3$, which is exactly the three free parameters
$(a,\alpha,\beta)$ of the period-one symbol --- a check I did not build in, and
the reason I believe the rest of the table.

So the reachable symbols form a $(2d+1)$-dimensional family, and every $G$ is
locally reachable precisely when $2d+1\ge k+1$, i.e.
\[
  d\ \ge\ \lceil k/2\rceil ,
\]
not $d\ge(k+2)/5$. For $k=6$ that means period \emph{three}, not period two ---
which is presumably why the period-two work of Entry 12 stalled, and why my
first attempts at $k=6$ with $d=2$ found nothing. I had been running the wrong
experiment for a reason I had written down incorrectly the same day.

Two things I want to keep straight about this. It is an \textbf{observation
supported by computation}, not a theorem: I have not proved the rank formula,
only verified it over a range. And it is a statement about \emph{local}
reachability --- a submersion means nearby symbols are reachable, not that a
specific target $G$ with all roots on the grid is. Realisability (real weights,
no vanishing edge) remains a separate hurdle, and it is exactly the hurdle that
killed the period-two construction and every rational candidate at $k=3$.

What makes it worth the entry is the shape of the consequence. The cyclotomic
certificates have moduli that are lcms of cyclotomic orders --- $60$, $70$,
$90$, $24$ --- which is why the results read as a list of congruence classes. A
period-$d$ construction has modulus $d$. If the count is the only obstruction,
the right statement is not ``$n$ divisible by $60$'' but ``$n$ divisible by
$\lceil k/2\rceil$, once $n/d$ is large enough to hold $k+1$ grid roots'' ---
one clean congruence for every $k$, instead of a table.

\textbf{A methodological note.} The rank computation took a few minutes to
write and settled a question I had been attacking by search for hours. The
search was asking ``can I find a certificate here?''; the rank was asking ``how
many dimensions of certificate exist here at all?'' The second question was
cheaper and strictly more informative, and I should have asked it first.

\section*{Entry 18: $k=6$, and what it cost to get there}
\textbf{Date:} September 26, 2026 \\
\textbf{Commands:} {\small\texttt{prescribe\_symbol.py},
\texttt{verify\_period\_d\_cert.py}}

$Z(P(60,6))=Z(P(72,6))=14$. These are the first exact values at $k=6$, a range
where Entry 16 showed the period-one construction \emph{provably} cannot
reach the ceiling.

The route was not more search. Once the rank law said the symbol map is a
submersion for $d\ge\lceil k/2\rceil$, the right move was to stop searching for
a certificate and start \emph{prescribing} one: choose $k+1$ interior grid
values, form the monic $G^\ast$ with those roots, and solve $G(w)=G^\ast$. That
is a well-posed system with a full-rank Jacobian, where minimising the $r$
smallest eigenvalues is a badly-posed one that leaves the optimiser room to
cheat.

And it cheated exactly as you would expect. Every unconstrained solve returned
a point with one spoke weight collapsed to zero --- a real solution of
$G(w)=G^\ast$, because with a spoke gone the matrix block-diagonalises and $G$
factors as $\det U\cdot\det V$, but a matrix for a \emph{different graph}.
Three things failed to fix it: a realisability barrier (which demonstrably
works --- it improved the known $k=5$ certificate's margin from $0.288$ to
$0.321$ without breaking it); continuation, which lost the symbol at the first
step $\tau=0.02$; and simply restarting, because these are \emph{singular}
points where the Jacobian rank drops from $7$ to $6$, and singular points are
attractors.

What fixed it was noticing that the failures were all at $n=48$ and $n=54$,
where $\Gamma_m$ has $7$ and $8$ interior values against the $7$ required ---
one admissible root set, then eight. At $n=60$ and $n=72$ there are $9$ and
$11$, hence $36$ and $330$ root sets, and realisable points exist. The
obstruction was never really the degeneracy; it was that I had been testing
the two cases with no room to move.

\textbf{A mistake worth recording.} I had queued $n=24,30,36,42$ for the first
$k=6$ run. Every one of them is impossible by construction: my guard required
$m\ge k+2$, but attaining $2k+2$ needs $k+1$ distinct \emph{interior} grid
values, so the requirement is $m\ge2k+4$. $\Gamma_8$ has three interior values
where seven are needed. Had those runs come back negative --- and they would
have --- I would have read it as evidence against period three rather than
evidence of a badly designed experiment.

\textbf{A mechanism I proposed and killed.} Seeing that the $n=48$ target had
constant term zero (because $4\mid16$ puts $\sigma=0$ in the grid, and with no
slack it had to be included), I proposed that a zero constant term forces the
spokes to vanish, and predicted $n=54$ would escape. It does not: $G(0)$ is
nonzero when the spokes are zeroed, and does not scale with them. I had
inferred a structural law from one data point, and the test took two minutes.

\textbf{Status, stated precisely.} These certificates are verified on the full
$2n\times2n$ matrix --- not the Fourier blocks that found them --- with fourteen
eigenvalues below $4\times10^{-15}$, the fifteenth at $5\times10^{-2}$, the
correct cubic pattern, and every required entry near a fifth of the matrix
scale. Several independent root sets give certificates at each $n$, so this is
a region and not a lucky point. But the weights are algebraic numbers from a
Newton solve, not integers, so this is \textbf{numerically certified and not
yet exact} --- a weaker standard than the $k=4$ and $k=5$ results, and I am not
going to describe it as though it were the same. Choosing the root set to be a
union of full Galois orbits makes $G^\ast$ rational and is the route to closing
that gap; it is running.

\section*{Entry 19: The method scales --- $k=7$ and $k=8$}
\textbf{Date:} September 26, 2026 \\
\textbf{Command:} {\small\texttt{python verification/prescribe\_symbol.py}}

$Z(P(96,7))=16$ and $Z(P(96,8))=18$, on top of $Z(P(60,6))=Z(P(72,6))=14$.

What makes this worth more than three extra values is that the period was
\emph{predicted}. The rank law says the symbol map becomes surjective at
$d\ge\lceil k/2\rceil$, so $k=6$ needs period three and $k=7,8$ need period
four. That is what happened: $d=3$ carried $k=6$ and failed for $k=7$ and
$k=8$, where $d=4$ succeeded. A law measured on one set of cases went on to
predict, correctly, which experiments would work before they were run. That is
the difference between a fitted pattern and a structural fact, and it is the
main reason I now believe the rank law rather than merely reporting it.

The claim also changes shape. Before today the results were a list ---
$k=4$ here, $k=5$ there, on scattered congruence classes. Now there is a
\emph{method} with an exact ceiling equal to the known upper bound, a rule
saying how much symmetry must be broken to reach it, and certificates
attaining it for every $k$ from $2$ to $8$.

\textbf{Generality.} A sweep over $n$ for $k=6$ at $d=3$ finds certificates at
$n=60$, $72$, $84$ --- the last two on the first or second root set tried. So
this is not two lucky values of $n$; once the counting conditions
($d\ge\lceil k/2\rceil$ and $m=n/d\ge2k+4$) hold with a little slack in the
choice of roots, certificates appear readily.

\textbf{What is still weaker than it looks.} All of $k\ge6$ is numerically
certified, not verified in exact arithmetic. The route I expected to fix that
--- choosing the root set to be a union of full Galois orbits, making the target
symbol rational --- is \textbf{closed}: all eight Galois-closed sets at $n=72$
hit the symbol to $10^{-15}$ with nullity $14$, and all eight collapse a
weight, whether or not $\sigma=0$ is among the roots. Rationality and
realisability are in tension here, and I do not understand why. The remaining
route is interval Newton, which the ceiling theorem makes tractable: certifying
the \emph{equation} $G(w)=G^\ast$ exactly gives the nullity as a discrete
consequence, since $k+1$ grid roots force $2k+2$ singular blocks and the
ceiling forces equality. That converts a fragile numerical rank claim into a
genuine computer-assisted proof.

\textbf{A process note.} Twice today the decisive step was to stop searching
and compute a dimension instead --- first the rank of the symbol map, then the
count of admissible root sets. Both took minutes and both redirected hours of
search. The failures at $n=48$ and $n=54$ looked like a structural obstruction
and were nothing of the kind: those are the two cases with one and eight
admissible root sets, against thirty-six and three hundred and thirty at
$n=60$ and $n=72$. I had been testing the only cases with no room to move.

\section*{Entry 20: Enlarging the object, and a classification}
\textbf{Date:} September 26, 2026 \\
\textbf{Commands:} {\small\texttt{cyclic\_covers.py},}\\
{\small\texttt{adjacency\_classification.py}}

Two changes of framing, both of which make the work about something larger than
$P(n,k)$ without weakening anything.

\textbf{Cyclic covers.} $P(n,k)$ is a cyclic cover: a two-vertex base graph
with loops of voltage $1$ and $k$ and a spoke of voltage $0$, lifted to
$V(B)\times\mathbb Z_n$. Nothing in the Fourier reduction ever used the
particular voltages, only that $\mathbb Z_n$ permutes the fibres --- so the
block reduction, and the ceiling as the degree span of $\det M(\zeta)$, hold
for \emph{any} base graph. Verified on seven base graphs on two to five
vertices. The right setting was covering-space theory all along; I had been
proving a special case without noticing.

\textbf{A classification, via Conway--Jones.} The adjacency nullity is
$\#\{m: 2\cos\frac{2\pi(k+1)m}{n}+2\cos\frac{2\pi(k-1)m}{n}=1\}$, a rational
relation among cosines of rational angles. Clearing denominators with
$z=\omega^m$ turns it into: which roots of
$P_k(z)=z^{2k+2}+z^{2k}-z^{k+1}+z^2+1$ are roots of unity? So
\[
  \operatorname{null}(\mathrm{Adj}\,P(n,k))=\sum_{\Phi_d\mid P_k,\ d\mid n}\varphi(d),
\]
exactly, for every $k$ and $n$. The adjacency matrix is \emph{never} singular
for $k=3,6,9$; it reaches the full ceiling only at $k=2$ and $k=5$. Checked
against a direct count at every $13\le n<200$ for $k=2,5,8$.

\textbf{An error this exposed.} The paper claimed ``in the computed range
$k\le10$ only $k=2,5,8$'' admit solutions. That came from a scan I had filtered
at nullity $\ge6$, which silently discarded $k=4,7,10,11,12$ --- all of which do
admit solutions, at nullity $2$ or $4$. The filter was a display choice that
became a mathematical claim. Corrected.

\textbf{A weak test I nearly reported as strong.} My first cyclic-cover check
compared the full matrix nullity against the sum of block nullities and passed
everywhere --- but every case had nullity $0$, so it only confirmed $0=0$. I had
to tune the weights to force genuinely singular blocks before the test meant
anything. It then passed, and exposed a second problem: one base graph had a
loop of voltage $0$, which lifts to a self-loop rather than an edge, so my two
builders disagreed. Both were my errors, not the theory's, but a passing test
that tests nothing is worse than a failing one.

\textbf{Why enlarge the object at all.} Partly because it is the truer
statement. But also, candidly, because the competitive field rewards it: a
result about cyclic covers of graphs reads as mathematics, while a result about
one parametrised family reads as a case study --- even when the theorem is
identical. The attributed open problem survives untouched as the worked
application, so nothing is traded away.

\section*{Entry 21: Four investigations, two negative}
\textbf{Date:} September 26, 2026

\textbf{1. Why exactness is hard for $k\ge6$ --- a real obstruction.} If the
period-$d$ weights are rational then the Galois automorphism
$\sigma_j:\zeta_m\mapsto\zeta_m^{\,j}$ carries $M_\ell$ to $M_{j\ell}$ and fixes
the weights, so $\sigma_j(\det M_\ell)=\det M_{j\ell}$ and the set of singular
blocks must be a union of full Galois orbits. A rational certificate can
therefore only realise a \emph{Galois-closed} root set --- and those are exactly
the sets that degenerate. So for $(k,d,n)=(6,3,60)$ and $(6,3,72)$ there is no
rational certificate at all. This makes the exactness question and the
Galois-closedness question the same question, which I had not seen.

Testing it taught me something about testing. My first check asked whether
$\det M_{j\ell}=\det M_\ell$ numerically. It failed everywhere, and for a
moment I believed the proposition was false. But those two numbers are
\emph{conjugates}, not equals --- the correct consequence is that the
\emph{zero set} is stable, equivalently that the norm over an orbit is
rational. Tested that way it holds exactly, to machine zero, in every case. The
proposition was right and my test was measuring the wrong thing.

\textbf{2. The rank law: half explained.} The nullspace of the Jacobian of
(weights)~$\mapsto$~(symbol) has dimension $3d-1$. Conjugation by a positive
diagonal matrix preserves the pattern, the nullity, and the monic symbol, and
accounts for exactly $2d$ of those --- verified directly. The remaining $d-1$ I
cannot identify. My hypothesis was that the symbol sees the off-diagonal
weights only through their products, as happens for the determinant of a
periodic Jacobi matrix; rescaling $b$, $c$ or $e$ at fixed product changes the
symbol in every case tested, so that is wrong. The rank law stays an
observation.

\textbf{3. The ceiling on general covers.} Sampling random matrices tests
nothing here, since they are generically nonsingular --- my first attempt made
exactly that mistake and passed trivially, confirming $0=0$. Maximising the
nullity instead, over fully non-equivariant weights on four small covers, the
largest attainable value never exceeded the degree span and equalled it twice.
Evidence, not proof.

\textbf{4. Ihara zeta functions --- checked, and not claimed.} The same Fourier
decomposition gives the full adjacency spectrum of $P(n,k)$ in closed form,
$\lambda_\pm=\tfrac{s+t}{2}\pm\sqrt{(\tfrac{s-t}{2})^2+1}$, confirmed against
direct diagonalisation to $10^{-15}$. From it one reads off that $P(n,k)$ is
Ramanujan for small $n$ but never for large $n$, since eigenvalues accumulate
at $3$. But that spectrum is already known and the conclusion is unsurprising,
so there is no contribution here and I am not claiming one. I record it because
the point of checking a speculative connection is to be willing to drop it.

\section*{Entry 22: A subfield ladder that does not (yet) hold weight}
\textbf{Date:} September 26, 2026 \\
\textbf{Command:} {\small\texttt{python verification/subfield\_certificates.py}}

Entry 21 proved that rational weights force a Galois-closed root set, and every
Galois-closed set at $k=6$ degenerates, so no rational certificate exists at
$(6,3,60)$ or $(6,3,72)$. I thought the way around it was to run the same
argument over a subfield: weights in $K=\mathbb Q(\zeta_m)^H$ need only an
$H$-stable singular set, $H$-orbits are smaller, so more root sets qualify.

At $(6,3,72)$ all three index-two subgroups do admit realisable root sets:
$H=\{1,7\}$ with $\ell=(1,2,3,6,7,9,10)$, $H=\{1,11\}$ with
$\ell=(1,2,3,6,8,9,11)$, $H=\{1,5\}$ with $\ell=(1,2,3,5,8,9,10)$, whose fixed
fields are $\mathbb Q(\sqrt2)$, $\mathbb Q(\sqrt3)$, $\mathbb Q(\sqrt6)$. I
wrote this up as ``the obstruction lifts''.

\textbf{That was a logical error, and I am recording it rather than quietly
fixing it.} The implication runs one way only:
\[
  \text{weights in } K \ \Longrightarrow\ \text{root set is } H\text{-stable}.
\]
The converse does not follow. An $H$-stable root set being realisable says only
that the \emph{necessary} condition is not violated; it gives no reason to
believe any certificate for that set has weights in $K$. I treated a necessary
condition as if it were sufficient.

The evidence since is against it. PSLQ on the extracted $H=\{1,7\}$ certificate
finds no relation of small height in $\mathbb Q(\sqrt2)$, and against a
degree-four basis returns coefficients of height $10^{2}$--$10^{3}$, which is
noise rather than a relation. Parametrising the weights \emph{inside}
$\mathbb Q(\sqrt2)$ from the start --- using the reduction of the conditions to
$H$-orbits, only five independent rather than seven --- found nothing in sixty
trials. At $(6,3,60)$ every index-two set degenerates outright and only the
full degree-four field admits a realisable set. And at $(6,4,64)$, the one
place where the root set is forced to be everything and is therefore
automatically Galois-closed, a realisable certificate exists but none of its
free coordinates is rational.

\textbf{Where that leaves exactness for $k\ge6$.} Open, and the evidence now
leans negative. What stands is the proposition of Entry 21: no \emph{rational}
certificate exists at the parameters tested, and that is a theorem. What does
not stand is any claim to have found a small field that works. The honest
summary is that the $k\ge6$ certificates appear to require at least the full
real cyclotomic field, and I have not succeeded in recognising them even there.

The machinery built along the way is sound and stays: Gauss periods give a
$\mathbb Q$-basis of a fixed field, integer-relation detection proposes exact
coefficients, and the blocks are rebuilt over $\mathbb Q(\zeta_m)$ with
$\det M_\ell$ checked to be exactly zero. Two bugs in it were found and fixed
--- the periods are not a basis when the Ramanujan sum vanishes ($\eta_1+\eta_7=0$
for $m=24$, $H=\{1,5\}$) or when they span fewer dimensions than the degree,
and my first recognition routine fitted one real number against a
$D$-dimensional basis by least squares and then truncated to double precision
when integer-relation detection needs sixty digits.

\textbf{The lesson I want to keep.} I have now twice in one day turned a
one-directional implication into a two-directional one in my head: once with
$\sigma=0$ and the vanishing constant term, and once here. Both times the
statement I wanted to be true was the converse of the statement I had proved.
The tell is the same in both cases --- a result that arrives without any new
work being done.

\section*{Entry 23: The result is about covers, not about $P(n,k)$}
\textbf{Date:} September 26, 2026 \\
\textbf{Commands:} {\small\texttt{cover\_certificates.py},
\texttt{cover\_ceiling\_sharp.py},}\\
{\small\texttt{cover\_zf.py}}

The construction is not about generalized Petersen graphs. Prescribing
$D/2$ interior grid values and solving for the weights attains the ceiling on:
the $P(n,2)$ base; the theta graph with voltages $(0,1,3)$, a genuinely
different cubic family; a three-vertex path base with $D=12$ whose cover has
$3n$ vertices; and a four-vertex star base with $D=12$. In each case the
nullity is exactly $D$ with every required entry a healthy fraction of the
matrix scale. With $M\le Z$ this gives $Z(B^n)\ge D$ for any cover where the
construction succeeds --- a statement about a class, not a family.

\textbf{A hypothesis I proposed and then watched fail to be necessary.} I had
worked out that the $P(n,k)$ elimination generalises when the base graph admits
an ordering along which each equation releases one new unknown --- that is, when
it has a Hamiltonian path --- and was ready to state that as the hypothesis. The
four-vertex star has no Hamiltonian path and the ceiling holds there anyway.
So the condition is sufficient for my particular argument and not necessary for
the result. Better to know that before writing it down as a theorem.

\textbf{A test that was worthless until it was not.} My first check of the
ceiling on covers asked whether $\ker A$ injects into a window of $D$
consecutive fibre coordinates. It passed everywhere --- with matrices of nullity
$2$ against windows of length $6$ to $12$, which tests nothing at all. The
check only bites at the ceiling, and getting there required building
certificates of nullity exactly $D$ first. Repeated properly, injectivity holds
in every case.

\textbf{And a claim of mine that is simply false.} I wrote down
$Z(B^n)\ge D$ unconditionally. Exhaustive computation gives $Z(P(8,2))=5$ while
$D=6$. The bound needs the construction to actually produce a certificate, and
at $n=8$ it cannot: $\Gamma_8$ has exactly three interior values for the three
required, the no-slack configuration that degenerates. What holds
unconditionally is $Z(B^n)\ge\nu$ for whatever nullity $\nu$ is attained, and
that I checked against exhaustively computed zero forcing numbers on covers of
$14$ to $20$ vertices with no violation.

The interesting wrinkle is that slack is not the criterion either. For the
$P(n,2)$ base the construction attains the ceiling at $n=7$ with \emph{no}
slack, and fails at $n=10$ with slack to spare. Whether it succeeds is an
arithmetic property of $\Gamma_n$, which is the same phenomenon that governed
the cyclotomic certificates and the rational obstruction. That keeps being the
real content of this problem, and I do not understand it.

\section*{Entry 24: the obstruction, in closed form}
\textbf{Date:} September 26, 2026 \\
\textbf{Commands:} {\small\texttt{antipodal\_obstruction.py}, \texttt{zf n 2}}

Entry 23 ended by admitting that whether the construction succeeds ``is an
arithmetic property of $\Gamma_n$, which I do not understand.'' I understand it
now, and it is not arithmetic at all.

For $k=2$ the symbol has three roots and no residual conditions, so
realisability is a closed-form function of the three prescribed grid values.
Matching coefficients gives $c^{2}=e_3-e_1e_2$, and then the elementary
identity $(s_1+s_2)(s_1+s_3)(s_2+s_3)=e_1e_2-e_3$ turns it into
\[
  c^{2}=-(s_1+s_2)(s_1+s_3)(s_2+s_3).
\]
A factorisation, so $c^{2}=0$ if and only if two roots are antipodal. That is a
proof, not a pattern: the thing I had been treating as mysterious arithmetic is
one line of symmetric-function algebra.

\textbf{It is a parity phenomenon, for every $k$.} $L_k(-s)=(-1)^{k}L_k(s)$,
so for even $k$ the two root conditions at $y$ and $-y$ add to
$\beta=-aL_k(y)$ and subtract to $\alpha=-L_k(y)$, giving $c^{2}=a\alpha-\beta=0$
identically --- for every even $k$, not just $k=2$. For odd $k$ the same
elimination gives $c^{2}=L_k(y)(y^{2}-a^{2})/y$, and there is no obstruction at
all. I checked this symbolically for $k=2,\dots,8$: zero for $2,4,6,8$, a
nonvanishing rational function for $3,5,7$.

\textbf{Two exceptions, and they are not the same kind of thing.}
$\Gamma^{\circ}_n$ has no antipodal pair when $n$ is odd, and $n/4$ or
$(n-2)/4$ classes when $n$ is even, so three pairwise non-antipodal values
exist exactly for odd $n\ge7$ and even $n\ge12$. The exceptions above $6$ are
$n=8$ and $n=10$, in both cases by pigeonhole. Then exhaustive computation:
$Z(P(8,2))=5$, so there the method is right to fail, and $Z(P(10,2))=6$, so
there it is the method that falls short. I had been reporting both as ``the
construction fails''; they are opposite situations.

\textbf{The non-symmetric class finally earns its place.} At $n=7$ the only
triple gives $c^{2}=-1$: no symmetric certificate exists. In the
combinatorially symmetric class the two spoke weights need only have product
$-1$, and
\[
  A=\begin{pmatrix} I+P+P^{-1} & I\\ -I & P^{2}+P^{-2}\end{pmatrix}
\]
is an integer matrix of nullity exactly $6$ for every $n$ divisible by $7$. Up
to now the $\mathrm{cs}$ generalisation had bought me nothing --- every value
it certified, the symmetric class certified too, and I said so. This is the
first case where it is necessary.

\textbf{And it explains the old table.} The exhaustive rational enumeration at
$k=2$ rejected $12$ of $17$ factor sets, and I had recorded that as twelve
separate arithmetic accidents. They split cleanly: seven have $c^{2}$ exactly
zero and every one of them is antipodal, and the other five have $c^{2}=-1$ and
none of them is antipodal --- so all five are realisable in the larger class.
Ten of seventeen, not five, and the antipodal obstruction is the whole of what
remains. Corollary: for even $k$ the factor set may not contain both $d$ and
$2d$ for odd $d$, nor any $d$ divisible by $4$ with $\varphi(d)\ge4$. Checked
against $50$ candidate sets at $k=2,4$, no violations.

What I want to flag for myself is how long I sat on this. The closed form is
three lines of algebra on a formula that has been in the paper since Entry~15,
and I spent Entries 19--23 doing numerical searches around it.

\section*{Entry 25: the obstruction was the wrong question}
\textbf{Date:} September 26, 2026 \\
\textbf{Commands:} {\small\texttt{cs\_exhaustive.py}, \texttt{k3\_certificate.py},}\\
{\small\texttt{signed\_adjacency.py}}

Entry 24 found that for even $k$ an antipodal pair of symbol roots forces
$c^{2}=0$, and I went looking for the odd-$k$ analogue. There isn't one, because
antipodality was never the right statement. The right one is a single line:
\[
  c^{2}=a\alpha-\beta=0
  \iff \beta=a\alpha
  \iff F=(s+a)L_k+\alpha(s+a)=(s+a)\bigl(L_k+\alpha\bigr).
\]
So $c^{2}=0$ iff the symbol factors as $(s+a)(L_k+t)$. Every one of the eight
vanishing cases at $k=3$ is of that form --- $\Psi_9=L_3+1$, $\Psi_{18}=L_3-1$,
$\Psi_4\Psi_{12}=L_3$ --- and the antipodal theorem is the special case that
parity forces when $k$ is even. Two entries spent on a corollary.

\textbf{And the degenerate list is finite, by Niven.} If $F$ is rational with
all roots on the grid then $a$ is its $s^{k}$ coefficient, hence rational, and
$-a$ is a grid value, so $a=2\cos(\text{rational}\cdot\pi)$ is rational; by
Niven's theorem $a\in\{0,\pm1,\pm2\}$. Same for $t=-L_k(r)=-2\cos(k\varphi)$.
At most $25$ degenerate rational symbols for every $k$, explicitly listable.
Enumerating the $25$ pairs gives $7,8,7,6,9$ for $k=2,\dots,6$, matching the
independent exhaustive count exactly in each case.

\textbf{What I had been getting wrong for weeks.} The old enumeration required
$c^{2}>0$, the condition for a \emph{symmetric} certificate, and recorded
everything else as ``not realisable.'' But $c^{2}$ is the \emph{product} of the
two spoke weights, and in the combinatorially symmetric class those are
independent. Only $c^{2}=0$ is fatal. Redoing the enumeration with $c^{2}\ne0$:

\begin{center}\small
\begin{tabular}{crrrrr}
\toprule
$k$ & candidates & outside family & $c^{2}=0$ & sym & cs \\
\midrule
2 & 17 & 0 & 7 & 5 & 5 \\
3 & 35 & 21 & 8 & 0 & \textbf{6} \\
4 & 67 & 54 & 7 & 3 & 3 \\
5 & 127 & 118 & 6 & 1 & 2 \\
6 & 225 & 216 & 9 & 0 & 0 \\
7 & 385 & 378 & 6 & 0 & \textbf{1} \\
8 & 651 & 644 & 7 & 0 & 0 \\
9 & 1065 & 1057 & 8 & 0 & 0 \\
10 & 1705 & 1698 & 7 & 0 & 0 \\
\bottomrule
\end{tabular}
\end{center}

\noindent
Two cases I had written down as \emph{impossible for rational symbols} are not.
Every non-fatal rejection has $c^{2}=-1$ or $-2$.

\textbf{$k=3$ closes.} $\Psi_5\Psi_{10}=s^{4}-3s^{2}+1=sL_3(s)+1$, so
$a=\alpha=0$, $\beta=1$, $c^{2}=-1$, and
\[
  A=\begin{pmatrix} P+P^{-1} & I\\ -I & P^{3}+P^{-3}\end{pmatrix}
\]
has nullity exactly $8$ for every $n$ divisible by $10$. So $Z(P(n,3))=8$,
where the paper had $\ge7$. I checked it against exhaustive $Z$ at $n=10,20$
and against the nullity at $n$ up to $100$. The matrix is the adjacency matrix
with one set of spokes negated, and the sign carries the whole result: unsigned,
the nullity is $0$.

Its four roots are $\pm\varphi,\pm\varphi^{-1}$ --- \emph{two antipodal pairs},
which for even $k$ would be fatal. So the parity corollary is not only an
obstruction, it says where to look when $k$ is odd. That is the first time in
this project a negative result has pointed at a positive one.

\textbf{$k=7$ closes too, and the two are the same phenomenon.} The $k=7$
symbol also has $a=\alpha=0$, $\beta=1$. Setting $s=z+z^{-1}$ turns
$sL_k+1=0$ into
\[
  Q_k(z)=z^{2k+2}+z^{2k}+z^{k+1}+z^{2}+1=0,
\]
which is the adjacency polynomial $P_k$ of Entry~18 with the middle sign
flipped. So the whole classification I proved for the adjacency matrix applies
verbatim to the signed one, and the question ``which $k$ admit an integer
certificate of this shape'' becomes ``for which $k$ does $Q_k$ split into
cyclotomics.'' Answer, for $2\le k\le60$: only $k=2,3,7$, giving
$\Phi_5\Phi_6$, $\Phi_5\Phi_{10}$, $\Phi_5\Phi_{10}\Phi_{24}$ and the ceiling
each time. Hence $Z(P(n,7))=16$ for every $n$ divisible by $120$ --- the paper
previously had $k=7$ only from a floating-point period-$4$ certificate at a
single $n$. The unsigned $P_k$ splits only at $k=2,5$, so the two sign choices
together cover $k\in\{2,3,5,7\}$ and nothing else.

I want to be clear that this is rigidity, not an infinite family. I checked to
$k=60$ hoping for more and there are none, which is the same Lehmer-type
behaviour the unsigned case showed.

\textbf{The negative side got stronger as well.} $k=6,8,9,10$ are impossible
over $\mathbb Q$ in \emph{both} classes, and no longer because a search came up
empty: every candidate that clears the family condition is one of the finitely
many degenerate symbols, and the Niven bound caps that list at $25$
before any computation. At $k=6$ all nine survivors are $(s+a)(L_6+t)$ with
$a,t\in\{0,\pm1\}$.

The lesson I should have learned three entries ago: when a computation returns
``not realisable,'' check which inequality was being tested.

\section*{Entry 26: off the congruence classes}
\textbf{Date:} September 27, 2026 \\
\textbf{Commands:} {\small\texttt{general\_n.py}, \texttt{krawczyk.py},
\texttt{certify\_pipeline.py},}\\
{\small\texttt{gauge\_fixed.py}, \texttt{certify\_sweep.py}, \texttt{tiles.py}}

Every exact value in this project has been attached to a divisibility class.
That is not a gap in the searching; it is forced, and I can now say how.

A period-one symbol is a \emph{fixed} polynomial. Its roots are fixed algebraic
numbers, and a fixed algebraic number lies on $\Gamma_n$ only for $n$ in a
divisibility class. So the set of $n$ a given period-one matrix can settle is
$\{n:\operatorname{lcm}(D)\mid n\}$, which has density at most $1/3$ and
contains \emph{at most one prime}, since $\{n:L\mid n\}$ with $L\ge3$ contains a
prime only when $L$ is that prime. To cover more, the weights must depend on
$n$.

For $k\ge3$ that dependence is constrained by one identity I had not noticed.
Since $L_k$ has nonzero coefficients only in degrees congruent to $k$ mod $2$,
the coefficient of $s^{k-1}$ in $(s+a)L_k+\alpha s+\beta$ is the $s^{k-2}$
coefficient of $L_k$, namely $-k$, for \emph{every} $(a,\alpha,\beta)$. So every
period-one symbol satisfies
\[
  e_2(r_1,\dots,r_{k+1})=-k .
\]
On the grid, after expanding $e_1^2$ by product-to-sum, that is a rational
linear relation among cosines of rational angles --- a Conway--Jones relation,
hence rigid. For $k=2$ the coefficient of $s^{k-1}=s$ is $\alpha-2$, free, so
the condition is vacuous. \emph{That} is why $k=2$ is general in $n$ and nothing
else is: the $k=2$ theorem gets to choose its roots as functions of $n$.
I checked $e_2=-k$ against all $26$ certificates in the paper; it holds in every
case with $k\ge3$ and takes six different values among the $k=2$ ones.

\textbf{So stop using the invariant matrices.} The ceiling corollary
characterises nullity $2k+2$ for \emph{any} matrix on the pattern by $T=I$ on
the monodromy, and that condition has no arithmetic in it. Minimising
$\|T-I\|$ over all $5n$ weights reaches the ceiling at $k=2$ for
$n=9,11,13,14,15$ and at $k=3$ for $n=17,20$ --- including primes.

\textbf{Certifying it, and the mistake I nearly shipped.} $T=I$ is bad to
certify: the image of $w\mapsto T$ is $51$-dimensional inside the
$64$-dimensional space of $8\times8$ matrices at $k=3$, so no square subsystem
is available. Certify the nullity instead, via $A(w)K(X)=0$ with $K$ in graph
form --- bilinear, so the Jacobian is affine and its Lipschitz constant is an
explicit integer, which makes a Krawczyk bound cheap.

I first chose the square subsystem by QR pivoting, got a passing Krawczyk test,
and was ready to call it a proof. It was not one: the $28$ equations I had
dropped were not shown to vanish at the true zero, only at my starting point.
The fix is structural. $K^{\!\top}\!AK$ is symmetric because $A$ is, and
$28=\binom82$ is exactly its antisymmetric part; if the non-pivot rows vanish
then $K^{\!\top}\!AK$ equals the pivot block, which must therefore be symmetric,
so setting its upper triangle to zero forces the lower triangle. The system
\emph{all non-pivot rows} $+$ \emph{upper triangle of the pivot block} is
equivalent to $AK=0$, not merely a maximal-rank selection of it. And the
reduction is checkable: the $28$ dependent equations fall to $10^{-61}$ without
being solved for.

With that, a proof:
\[
  Z(P(n,3))=M(P(n,3))=8\qquad n=17,19,21,23,25,27,
\]
the primes $17,19,23$ among them. Combined with the $10\mid n$ certificate this
holds at $n=10,17,19,20,21,23,25,27,30,40,\dots$, a set no divisibility
condition describes. First results in this project not living on a congruence
class. The certificate at $n=17$ was found three independent ways and all three
certify; since stage~1 is only a search and stages~2--3 are the proof, the route
does not enter the result.

\textbf{Three of my own claims corrected, one in the middle of relying on it.}
The threshold: I wrote $n\ge(2k+2)^2/5$ from a loose codimension count. Spending
the gauge properly --- the $2n$ diagonal-conjugation parameters normalise every
spoke to $1$ and, for odd $n$, every outer weight to $\pm1$ --- leaves $3n$
unknowns and gives $n\ge(k+1)(2k+3)/3$, which is $35$ at $k=6$, not $27$. I had
already run $k=6$ at $n=27,29,31$ and read their failure as meaningful. It was
not; those $n$ are below the threshold and could not have worked.

The growth rate: I conjectured $N(k)=O(k)$ from the saturation data
($\approx4k+3$). The dimension count is a necessary condition and forces
$N(k)$ to be at least quadratic. Withdrawn.

The monodromy is not symplectic. I guessed it would be, since the recurrence is
self-adjoint. The standard form is not preserved; the space of skew $B$ with
$T^{\!\top}BT=B$ has dimension $k+1$, so the eigenvalues come in reciprocal
pairs and the characteristic polynomial is palindromic. Different statement,
and it does not give the equation count I wanted from it.

\textbf{And one instrument withdrawn.} The eigenvalue-minimising search whose
plateaus I have been quoting as evidence against reachability gives maximum
nullity $7$ at $P(20,3)$, where $8$ is attained --- by the $\Psi_5\Psi_{10}$
certificate and, in the symmetric class, by the new method. It collapses onto
the degenerate strata, all spokes zero or all edges zero. So the conclusion I
drew from it, that exactness for $k\ge6$ looked unlikely, was never supported by
that data. I have marked every negative result from it as unreliable in the
paper rather than quietly dropping them.

\textbf{Toward all large $n$.} If a weight block has transfer product $I$, and
blocks compose, then $n$ is any sum of available block lengths. Blocks do not
compose naively, because $T_i$ depends on a window of width about $2k+2$; giving
every tile a fixed docking segment of $k+1$ positions at each end fixes that,
and I checked it rather than assumed it --- perturbing one tile's interior leaves
the other tiles' products bit-identical, $0.000\mathrm{e}{+}00$, for $k=2,3,4$.
Two coprime lengths would only give $n=ia+jb$ and leave everything below
$(a-1)(b-1)\approx600$ uncovered, which is useless; tiles of \emph{every} length
in $[L,2L)$ instead make every $n\ge L$ a sum, with no gap. For $k=3$ a tile of
length $\ell$ has $5(\ell-8)$ free weights against $64$ equations, so $\ell\ge21$
and $L=21$ would give all $n\ge21$. Whether those systems are solvable with all
weights nonzero is being computed and I do not know the answer yet.

\textbf{A regression I caused and caught.} Warm-starting each $n$ from the
previous certified solution made the sweep affordable, but I gated the cold
fallback on the warm start being worse than $10^{-2}$. A warm start returning
$10^{-3}$ then suppressed the cold solve and failed stage~2, and the sweep
reported only $n=17,19$ --- losing $21,23,25,27$, which I had already certified.
Restructured so stage~1 produces candidates and stage~2 is attempted on each in
turn. Worth recording because the failure mode was silent: the sweep looked like
a clean negative result.

\begin{thebibliography}{9}
\bibitem{RST} S. Rashidi, N. Shajareh Poursalavati, M. Tavakkoli,
\emph{J. Algebra Combin. Discrete Struct. Appl.} 7(2) (2020), 183--193.
\bibitem{K} A. Krishnan, arXiv:2607.19412 (2026).
\end{thebibliography}
\end{document}
